\section{\iflanguage{english}{Satzilla Outcome}{Satzilla Ergebnis}}\label{sec:satzilla_outcome}

\subsection{Validierung mit Fake Solvern}

Da ich in meinem ursprünglichen Bachelorprojekt Predictions in Satzilla hatte die nicht mit den 
tatsächlichen vorhersagen übereinstimmten (so wurde fälschlicherweiße ISAsat alle der beste Solver 
für viele Probleme empfohlen) riet mir mein Tutor Dr. Mathias Fleury zu testen ob Satzilla überhaupt
richtig funktioniert und passende vorhersagen machen kann.

Bevor ich Satzilla auf die Laufzeitdaten realer Solver angewandt haben, habe ich zunächst die gesamte 
Pipeline anhand künstlich erzeugter Solver-Verhaltensmuster validiert. Dieser Schritt war 
entscheidend, um sicherzustellen, dass die Merkmalsextraktion, die Datenausrichtung und der 
Lernmechanismus korrekt funktionieren.




\subsubsection{Design der fake Solver}




\subsubsection{Resultate}




\subsubsection{Interpretation}









\subsection{Trainieren auf echten Kissat Solvern}






\subsubsection{Dataset}









 \subsubsection{Training Konfiguration}






\subsubsection{Testgenauigkeit}




\subsection{Top-K-Genauigkeit: Der Wert des „fast Richtigen“}








\subsubsection{Resultate}






\subsubsection{Interpretation}






\subsubsection{Kumulative Genauigkeitskurve}








\subsection{Feature Importance Analysis}





\subsubsection{Wichtigste Features}






\subsubsection{Konklusion}







 \subsection{Per-Solver-Leistung}





 \subsubsection{Interpretation}





 \subsection{Vergleich mit einer zufälligen Referenzlinie}





 \subsection{Praktischer Nutzen für die SAT-Lösung}






 \subsection{Auswirkungen der Hinzufügung von fehenden LP-Features}






 \subsubsection{Interpretation}




 \subsection{Zusammenfassung der wichtigsten Ergebnisse}





 \subsection{Konklusion}

